\section{Place four complete groups in byte lanes}
An MSA register contains sixteen bytes. Four input groups of three bytes occupy twelve
of them, and their sixteen sextets fill the output register. Keep the group arithmetic
identical to the scalar calculation; only the selected lane positions change.

The encoder kernel first copies exactly twelve source bytes into sixteen private bytes,
with the remaining four initialized to zero. This avoids a sixteen-byte caller load
for an input contract that promises only twelve bytes. A larger mapped arena does not
extend the selected readable interval. The subsequent \texttt{ld.b} reads all sixteen
private bytes. Byte-form MSA transfers do not require a sixteen-byte-aligned caller
address, but they still transfer sixteen bytes.

Two gathers provide the source bytes for the two contributions to each sextet. For the
first three-byte group, their selector and shift patterns are:
\tagpdfsetup{table/header-rows={1}}
\begin{center}
  \begin{tabular}{lrrrr}
    Output sextet position   & 0 & 1  & 2 & 3 \\
    \hline
    Left source index        & 0 & 0  & 1 & 2 \\
    Left shift               & 0 & 4  & 2 & 0 \\
    Then logical right shift & 2 & 0  & 0 & 0 \\
    Right source index       & 0 & 1  & 2 & 0 \\
    Right logical shift      & 0 & 4  & 6 & 0 \\
    Right mask               & 0 & 15 & 3 & 0 \\
  \end{tabular}
\end{center}

The next groups add three, six and nine to the useful source indices. A masked-out
contribution can select an ordinary existing byte because its mask makes that value
irrelevant. After shifting, combine the two contributions and mask every byte with 63.
Byte shifts truncate each lane independently, so bits cannot carry into the next sextet.
